% Vietnamese translation for OI Public Library
% Source-File: IOI中国国家候选队论文/2005/胡伟栋/附件/传染病控制解题报告.doc
\documentclass[11pt]{article}
\usepackage{oipl-vn}

\title{Lời giải bài ``Khống chế dịch bệnh''}
\author{}
\date{}

\begin{document}
\maketitle

\origtitle{Epidemic Control solution report}
\sourcepath{IOI 2005 / Hu Weidong / attachments / Epidemic Control solution report.doc}

\section{Tóm tắt đề bài}

Cho một cây, trong đó nút gốc \(1\) là bệnh nhân đã bị nhiễm. Trong mỗi chu kỳ
lây lan, chỉ được cắt một cạnh của cây; sau đó mọi nút con của các nút đang bị
nhiễm cũng sẽ bị nhiễm. Cần tìm một phương án cắt cạnh tốt nhất sao cho tổng
số nút bị nhiễm là nhỏ nhất.

\section{Phân tích lời giải}

Gọi tập các nút bị nhiễm có độ sâu lớn nhất là \term{nguồn lây mới nhất}, ký
hiệu \(S\). Sau \(x\) chu kỳ lây lan, độ sâu của nguồn lây mới nhất là cố
định, bằng \(x\), nếu lấy độ sâu của gốc \(1\) là \(0\).

Từ đó có thể nghĩ đến một quy hoạch động theo tập:
\[
  f(x,S)=\text{số nút bị nhiễm nhỏ nhất sau \(x\) chu kỳ, khi nguồn lây mới
  nhất là \(S\)}.
\]
Từ \(f(x,S)\), nếu \(C(S)\) là tập tất cả các nút con của các nút trong \(S\),
và \(v\in C(S)\) là nút ở đầu sâu hơn của cạnh bị cắt, thì trạng thái kế tiếp
có dạng
\[
  f\bigl(x+1,C(S)\setminus\{v\}\bigr).
\]
Cách này nhìn có vẻ trực tiếp, nhưng số trạng thái theo tập là quá lớn; độ
phức tạp bộ nhớ tăng theo cấp mũ của số nút trên một tầng, hoàn toàn không
thể cài đặt.

Vì khó khăn về bộ nhớ, ta chuyển sang tìm kiếm. Xét không gian trạng thái:
đề bài cho một cây. Gọi cây của đề bài là cây gốc; số tầng của cây tìm kiếm
chính là độ sâu của cây gốc, còn số nhánh của cây tìm kiếm liên quan đến độ
rộng của cây gốc. Nếu cây gốc rất sâu thì thường độ rộng nhỏ; nếu cây gốc rất
rộng thì độ sâu lại nhỏ. Vì vậy quy mô tìm kiếm không lớn như cảm giác ban
đầu.

\section{Cách tìm kiếm}

Trong \(x\) chu kỳ, bệnh luôn lan đến đúng độ sâu \(x\). Nếu cắt một cạnh ở
độ sâu nhỏ hơn \(x\), bệnh đã đi qua cạnh đó nên việc cắt không còn ý nghĩa.
Nếu cắt một cạnh ở độ sâu lớn hơn \(x\), hiệu quả chắc chắn không bằng cắt
cạnh nối đến tổ tiên tương ứng ở gần nguồn lây hơn. Do đó trong một phương án
tối ưu, cạnh bị cắt ở chu kỳ hiện tại phải là một trong các cạnh nối từ nguồn
lây mới nhất \(S\) đến các nút con của nó.

Tìm kiếm như vậy đã vượt qua được phần lớn dữ liệu, nhưng với dữ liệu cực hạn
vẫn cần tối ưu thêm.

\section{Các tối ưu}

\paragraph{Tối ưu 1.}
Khi mới nhìn bài này, rất dễ nghĩ đến tham lam: mỗi lần chọn xóa cây con có
nhiều nút nhất. Tuy cách tham lam này có phản ví dụ, nó vẫn là một hướng thử
có xác suất cao dẫn tới nghiệm tốt sớm. Vì vậy khi duyệt các lựa chọn cắt
cạnh, ta sắp các cây con theo kích thước giảm dần.

\paragraph{Tối ưu 2.}
Dùng cắt tỉa tối ưu. Giả sử hiện đang tìm kiếm đến tầng \(d\), đã có
\(\mathit{sick}\) nút bị nhiễm. Gọi \(\mathit{tmp}\) là tổng số nút con của
các nguồn lây ở tầng \(d\). Vì trong một chu kỳ chỉ cắt được một cạnh, ở chu
kỳ kế tiếp sẽ có ít nhất \(\mathit{tmp}-1\) nút mới bị nhiễm. Do đó nếu
\[
  \mathit{sick}+\mathit{tmp}-1\ge \mathit{ans},
\]
ta có thể cắt nhánh tìm kiếm hiện tại.

\paragraph{Tối ưu 3.}
Khi tìm kiếm đến một hoặc hai tầng cuối, cấu trúc các cây con còn lại là như
nhau đối với mục tiêu, nên chọn tùy ý một nhánh để xóa rồi thoát khỏi vòng
lặp.

Sau khi thêm ba tối ưu trên, hiệu suất tăng lên rõ rệt. Dữ liệu chính thức lớn
nhất được giải trong khoảng \(0.2\) giây theo báo cáo nguồn.

\section{Chương trình tham khảo}

Chương trình Pascal dưới đây được giữ nguyên từ nguồn gốc.

\begin{verbatim}
{$R-,Q-,S-,I-}
{$M 65521,0,655360}
{Author:Zhou Gelin}
{Date:2005.1.11}
{QQ:379688236}
{MSN:zhougelin@hotmail.com}
{Test Report}
{P4 1.7G}
{Windows Xp SP1}
{Turbo Pascal 7.0}
{
EPIDEMIC.in1 = 10.0 (0.06s)
EPIDEMIC.in2 = 10.0 (0.08s)
EPIDEMIC.in3 = 10.0 (0.05s)
EPIDEMIC.in4 = 10.0 (0.03s)
EPIDEMIC.in5 = 10.0 (0.03s)
EPIDEMIC.in6 = 10.0 (0.03s)
EPIDEMIC.in7 = 10.0 (0.06s)
EPIDEMIC.in8 = 10.0 (0.03s)
EPIDEMIC.in9 = 10.0 (0.20s)
EPIDEMIC.in0 = 10.0 (0.11s)
}
program epidemic;
const inputfilename='epidemic.in';
      outputfilename='epidemic.out';
      maxn=300;
type pnode=^node;
     node=record
       data:integer;
       next:pnode;
         end;
     list=array[1..maxn]of integer;
var n,p,maxdep,ans:integer;
    g:array[1..maxn]of pnode;
    father:array[1..maxn]of integer;
    dep:array[1..maxn]of integer;
    sum:array[1..maxn]of integer;
    deg:array[1..maxn]of integer;
    son:array[1..maxn]of ^list;
    a:array[1..maxn]of boolean;
procedure insert(a,b:integer);
var q:pnode;
begin
      new(q);
      q^.data:=b;
      q^.next:=g[a];
      g[a]:=q;
end;
procedure read_data;
var i,a,b:integer;
begin
      fillchar(deg,sizeof(deg),0);
      assign(input,inputfilename);
      reset(input);
      readln(n,p);
      for i:=1 to p do
      begin
         readln(a,b);
         insert(a,b);
         insert(b,a);
         deg[a]:=deg[a]+1;
         deg[b]:=deg[b]+1;
      end;
      close(input);
end;
procedure dfs(v:integer);
var q:pnode;
    i,j,k:integer;
begin
      if father[v]>0 then
      begin
        deg[v]:=deg[v]-1;
        dep[v]:=dep[father[v]]+1;
        if dep[v]>maxdep then maxdep:=dep[v];
      end
      else dep[v]:=1;
      getmem(son[v],2*deg[v]);
      sum[v]:=1;q:=g[v];i:=0;
      while q<>nil do
      begin
        if q^.data<>father[v] then
        begin
          i:=i+1;
          son[v]^[i]:=q^.data;
          father[q^.data]:=v;
          dfs(q^.data);
          sum[v]:=sum[v]+sum[q^.data];
        end;
        q:=q^.next;
      end;
      for i:=1 to deg[v]-1 do
      begin
         k:=i;
         for j:=i+1 to deg[v] do
            if sum[son[v]^[j]]>sum[son[v]^[k]] then k:=j;
         j:=son[v]^[i];
         son[v]^[i]:=son[v]^[k];
         son[v]^[k]:=j;
      end;
end;
procedure search(deep,sick:integer);
var i,j,tmp:integer;
begin
      if deep=maxdep then
      begin
        if ans>sick then ans:=sick;
        exit;
      end;
      tmp:=0;
      for i:=1 to n do
         if a[i] and (dep[i]=deep) then
           tmp:=tmp+deg[i];
      if tmp=0 then begin search(maxdep,sick); exit; end;
      if sick+tmp-1>=ans then exit;
      for i:=1 to n do
         if a[i] and (dep[i]=deep) then
           for j:=1 to deg[i] do
              a[son[i]^[j]]:=true;
      for i:=1 to n do
         if a[i] and (dep[i]=deep) then
           for j:=1 to deg[i] do
           begin
              a[son[i]^[j]]:=false;
              search(deep+1,sick+tmp-1);
              a[son[i]^[j]]:=true;
              if deep=maxdep-1 then break;
           end;
      for i:=1 to n do
         if a[i] and (dep[i]=deep) then
           for j:=1 to deg[i] do
              a[son[i]^[j]]:=false;
end;
procedure answer;
begin
      assign(output,outputfilename);
      rewrite(output);
      writeln(ans);
      close(output);
end;
var time:longint;
begin
      time:=meml[64:108];
      read_data;
      dfs(1);
      ans:=1000;
      fillchar(a,sizeof(a),0);
      a[1]:=true;
      search(1,1);
      time:=meml[64:108]-time;
      writeln(time);
      answer;
end.
\end{verbatim}

\end{document}
